\documentclass[a4paper,11pt]{ctexart}
\usepackage[margin=2cm]{geometry}
\usepackage{tikz}
\usetikzlibrary{positioning, arrows.meta, shapes.geometric}
\usepackage{booktabs}
\usepackage{tabularx}
\usepackage{xcolor}
\usepackage{tcolorbox}
\tcbuselibrary{skins,breakable}

% 配色定义
\definecolor{primary}{RGB}{24, 76, 120}
\definecolor{success}{RGB}{34, 139, 34}
\definecolor{danger}{RGB}{178, 34, 34}
\definecolor{boxbg}{RGB}{248, 250, 252}
\definecolor{codebg}{RGB}{240, 244, 248}

\newtcolorbox{detailcard}[2][]{
    enhanced, breakable,
    colback=boxbg, colframe=#2,
    fonttitle=\bfseries\small, title={#1},
    boxrule=0.8pt, rounded corners=3pt,
    top=6pt, bottom=6pt, left=8pt, right=8pt
}

\begin{document}

\begin{center}
    {\LARGE \textbf{大麦高并发抢票系统：订单生命周期与代码级状态机}} \\[6pt]
    {\small \textbf{Code-Level Finite State Machine (FSM) \& Concurrency Control Specification}}
\end{center}

\vspace{0.3cm}

% ==================== 1. 状态跃迁流向图 ====================
\begin{center}
\begin{tikzpicture}[
    ->, >=Stealth, thick,
    node distance=2.2cm and 5cm,
    % 状态节点样式
    state/.style={
        draw=primary, fill=primary!8, rounded corners=6pt,
        minimum width=3.8cm, minimum height=1.3cm,
        align=center, font=\small\bfseries
    },
    terminal/.style={
        draw=danger, fill=danger!8, rounded corners=6pt,
        minimum width=3.8cm, minimum height=1.3cm,
        align=center, font=\small\bfseries, dashed
    },
    paid/.style={
        draw=success, fill=success!8, rounded corners=6pt,
        minimum width=3.8cm, minimum height=1.3cm,
        align=center, font=\small\bfseries
    },
    trans/.style={
        font=\footnotesize\bfseries, align=center
    }
]

    % 节点定义
    \node[state] (CREATED) {
        0 : CREATED (待支付) \\
        \footnotesize \normalfont 初始态 $\cdot$ 扣减库存
    };

    \node[paid, right=of CREATED] (PAID) {
        1 : PAID (已支付) \\
        \footnotesize \normalfont 终态 $\cdot$ 锁定票权
    };

    \node[terminal, below=of CREATED] (CANCELLED) {
        4 : CANCELLED (已取消) \\
        \footnotesize \normalfont 终态 $\cdot$ 回滚自愈库存
    };

    % 跃迁连线
    \draw[->, line width=1.5pt, draw=success] (CREATED) -- node[trans, above=2pt, text=success] {
        【跃迁 A】用户付款成功 \\
        \scriptsize \texttt{payOrder(orderId)}
    } (PAID);

    \draw[->, line width=1.5pt, draw=danger] (CREATED) -- node[trans, right=4pt, text=danger] {
        【跃迁 B】超时关单 / 主动取消 \\
        \scriptsize \texttt{cancelOrder(orderId)}
    } (CANCELLED);

    \draw[->, line width=0.8pt, draw=gray, dashed] (PAID) -- node[trans, above=2pt, text=gray] {
        \scriptsize 后续生命周期 \\
        \scriptsize 2:SHIPPED / 5:REFUNDED
    } ++(3.2cm, 0);

\end{tikzpicture}
\end{center}

\vspace{0.2cm}

% ==================== 2. 代码级跃迁详情卡片 ====================

\begin{detailcard}[【跃迁 A 剖析】待支付 (0) $\rightarrow$ 已支付 (1) ：乐观锁与版本递增]{success}
\begin{itemize}[leftmargin=0.8cm, itemsep=3pt]
    \item \textbf{触发入口：} \texttt{OrderController.payOrder(@PathVariable Long id)}
    \item \textbf{第一道防线（状态门卫）：}
    \begin{quote}
        \footnotesize
        \texttt{OrderStatus current = OrderStatus.fromCode(order.getStatus());} \\
        \texttt{if (!current.canTransitionTo(OrderStatus.PAID)) \{} \\
        \quad \texttt{throw new BusinessException(400, "订单状态有误，无法支付！");} \\
        \texttt{\}}
    \end{quote}
    \item \textbf{底层执行 SQL（PostgreSQL 物理行锁与 CAS）：}
    \begin{quote}
        \footnotesize
        \texttt{UPDATE d\_ticket\_order} \\
        \texttt{SET status = 1, version = version + 1, pay\_time = NOW()} \\
        \texttt{WHERE id = ? AND version = 1;}
    \end{quote}
    \item \textbf{并发冲突拦截：} \texttt{int rows = mapper.updateById(order);} 若 \texttt{rows == 0}，说明在读取与更新之间已被关单线程修改，立即抛出 \texttt{BusinessException(409, "订单状态已发生变更")} 阻断脏更新。
\end{itemize}
\end{detailcard}

\vspace{0.2cm}

\begin{detailcard}[【跃迁 B 剖析】待支付 (0) $\rightarrow$ 已取消 (4) ：自动关单与库存自愈]{danger}
\begin{itemize}[leftmargin=0.8cm, itemsep=3pt]
    \item \textbf{触发入口：}
    \begin{itemize}
        \item 定时巡逻任务：\texttt{OrderTimeoutTask.scanAndCancelExpiredOrders()} (\texttt{@Scheduled(fixedDelay = 10000)})
        \item 用户主动取消：\texttt{OrderController.cancelOrder(@PathVariable Long id)}
    \end{itemize}
    \item \textbf{关单更新 SQL：}
    \begin{quote}
        \footnotesize
        \texttt{UPDATE d\_ticket\_order SET status = 4, version = version + 1 WHERE id = ? AND version = 1;}
    \end{quote}
    \item \textbf{两步级联库存自愈（\texttt{restoreStock}）：}
    \begin{enumerate}
        \item \textbf{数据库库存回滚：} \texttt{UPDATE d\_ticket\_category SET remain\_stock = remain\_stock + 1 WHERE id = ?;}
        \item \textbf{Redis 缓存一致性淘汰：} \texttt{redisTemplate.delete("ticket:category:" + id);}
    \end{enumerate}
\end{itemize}
\end{detailcard}

\vspace{0.2cm}

\begin{detailcard}[【并发竞态裁决】第 14 分 59 秒碰撞的底层原理]{primary}
\small
\textbf{冲突场景：} 用户在倒计时即将结束时点击支付，与后台定时关单线程在同一毫秒内到达。\\
\textbf{裁决机制：}
两个线程同时读取到 \texttt{version = 1}。当它们尝试执行 \texttt{UPDATE} 时，PostgreSQL 对该数据行施加**行级排他排队（X-Lock）**：
\begin{enumerate}[leftmargin=0.8cm, itemsep=2pt]
    \item 优先获取到行锁的线程执行成功，行内 \texttt{version} 跃迁为 \texttt{2}，返回影响行数 $1$。
    \item 后释放锁的线程执行匹配 \texttt{WHERE version = 1} 时发现版本已被改写，返回影响行数 $0$。
    \item \textbf{结论：} 永远只有一方能够胜出，绝对不会出现“钱扣了但票被关单”或“库存重复归还”的数据灾难。
\end{enumerate}
\end{detailcard}

\end{document}
